% =============================================================================
% Aula 03 — CRUD · Formulários · Projeto StockSales
% Beamer + tema Madrid | Curso Entra21 — Spring Framework
% Compilar: tectonic -X compile slides-aula-03.tex
% =============================================================================
\documentclass[aspectratio=169,11pt]{beamer}

\usetheme{Madrid}
\usecolortheme{default}

\usepackage[T1]{fontenc}
\usepackage[utf8]{inputenc}
\usepackage[brazil]{babel}
\usepackage{booktabs}
\usepackage{graphicx}
\usepackage{hyperref}
\usepackage{listings}
\usepackage{xcolor}
\usepackage{tikz}
\usetikzlibrary{arrows.meta, positioning, shapes.geometric}

\definecolor{codebg}{RGB}{245,245,245}
\definecolor{codegreen}{RGB}{40,120,40}
\definecolor{codeblue}{RGB}{30,70,160}
\definecolor{codered}{RGB}{180,50,50}

\lstset{
  basicstyle=\ttfamily\small,
  backgroundcolor=\color{codebg},
  frame=single,
  rulecolor=\color{black!20},
  breaklines=true,
  columns=fullflexible,
  keepspaces=true,
  showstringspaces=false,
  keywordstyle=\color{codeblue}\bfseries,
  commentstyle=\color{codegreen},
  stringstyle=\color{codered},
  tabsize=2,
  literate={á}{{\'a}}1 {é}{{\'e}}1 {í}{{\'i}}1 {ó}{{\'o}}1 {ú}{{\'u}}1
           {Á}{{\'A}}1 {É}{{\'E}}1 {Í}{{\'I}}1 {Ó}{{\'O}}1 {Ú}{{\'U}}1
           {ã}{{\~a}}1 {õ}{{\~o}}1 {Ã}{{\~A}}1 {Õ}{{\~O}}1
           {ç}{{\c{c}}}1 {Ç}{{\c{C}}}1
}

\newcommand{\onde}[1]{%
  \begin{block}{Onde estamos}
    #1
  \end{block}
}

\title{Aula 03 --- Spring Framework}
\subtitle{CRUD · Formulários · StockSales}
\author{Mauricio Duailibi Neto}
\institute{Turma Java Entra21}
\date{}

\begin{document}

\begin{frame}
  \titlepage
\end{frame}

\begin{frame}{Como será a aula de hoje}
  \begin{block}{Parte 1 --- 18h30 às 20h}
    \begin{itemize}
      \item Revisão rápida da Aula 02
      \item Como funciona uma requisição HTTP
      \item CRUD: cadastrar, listar, editar e excluir
      \item Projeto novo: \textbf{cadastro-alunos} --- fazemos juntos
      \item Desafios no mesmo projeto
    \end{itemize}
  \end{block}

  \vspace{0.4em}
  \begin{block}{Parte 2 --- 20h20 às 22h}
    \begin{itemize}
      \item Voltamos ao \textbf{StockSales}
      \item CRUD de produtos e de clientes
    \end{itemize}
  \end{block}
\end{frame}

\begin{frame}{Dois projetos nesta aula}
  Não vamos misturar os códigos.

  \vspace{0.6em}
  \begin{description}
    \item[cadastro-alunos] Projeto \textbf{novo}, só desta primeira parte.
      Aqui aprendemos CRUD com um cadastro simples de alunos.
    \item[StockSales] O sistema da loja, que já existe.
      Usamos na segunda parte, depois do intervalo.
  \end{description}

  \vspace{0.6em}
  Quando aparecer um código, o slide vai dizer:\\
  \textbf{qual projeto} e \textbf{qual arquivo} estamos mexendo.
\end{frame}

\begin{frame}{Mapa da aula}
  \tableofcontents
\end{frame}

% #############################################################################
\section{Revisão da Aula 02}

\begin{frame}{O que já sabemos}
  \begin{itemize}
    \item Criar um projeto no \texttt{start.spring.io}
    \item Página HTML fica em \texttt{src/main/resources/static/}
    \item Dados vêm de um Controller, em JSON
    \item O JavaScript busca esses dados com \texttt{axios.get}
  \end{itemize}

  \vspace{0.7em}
  No StockSales, a tela de produtos \textbf{só lista}.\\
  Ainda não dá para cadastrar, alterar nem apagar.
\end{frame}

\begin{frame}{Página e API}
  O Spring atende dois tipos de pedido:

  \vspace{0.5em}
  \begin{itemize}
    \item \textbf{Página:} \texttt{/} ou \texttt{/produtos.html}\\
          --- o arquivo HTML da pasta \texttt{static/}
    \item \textbf{API:} \texttt{/api/produtos}\\
          --- o método Java do Controller, devolvendo JSON
  \end{itemize}

  \vspace{0.6em}
  A pessoa vê a página.\\
  A página pede os dados para a API.
\end{frame}

% #############################################################################
\section{O que é uma requisição}

\begin{frame}{Navegador e servidor conversam}
  \begin{center}
  \begin{tikzpicture}[
      box/.style={draw, rounded corners, align=center, minimum width=3.2cm,
                  minimum height=1.2cm, font=\small\bfseries},
      seta/.style={-{Latex}, thick}
    ]
    \node[box, fill=blue!15] (nav) {Navegador\\{\normalfont\small a tela}};
    \node[box, fill=orange!20, right=3.2cm of nav] (srv) {Spring\\{\normalfont\small o servidor}};
    \draw[seta] ([yshift=4pt]nav.east) -- node[above, font=\small]{pedido} ([yshift=4pt]srv.west);
    \draw[seta] ([yshift=-4pt]srv.west) -- node[below, font=\small]{resposta} ([yshift=-4pt]nav.east);
  \end{tikzpicture}
  \end{center}

  \vspace{0.9em}
  Cada clique que busca ou grava dados faz um \textbf{pedido}.\\
  O servidor devolve uma \textbf{resposta}.

  \vspace{0.4em}
  Esse pedido se chama \textbf{requisição HTTP}.
\end{frame}

\begin{frame}{Uma requisição tem três partes principais}
  Pensem numa carta:

  \vspace{0.5em}
  \begin{enumerate}
    \item \textbf{O que você quer fazer} --- o verbo:\\
          GET, POST, PUT ou DELETE
    \item \textbf{Para onde} --- o endereço:\\
          \texttt{http://localhost:8080/api/alunos}
    \item \textbf{O conteúdo} --- o \textbf{corpo} da carta, se houver:\\
          os dados em JSON
  \end{enumerate}

  \vspace{0.6em}
  GET quase nunca leva conteúdo.\\
  POST e PUT quase sempre levam.
\end{frame}

\begin{frame}{GET: só pergunta}
  Exemplo: listar alunos.

  \vspace{0.5em}
  \begin{block}{Pedido que o navegador manda}
    \begin{itemize}
      \item Verbo: \texttt{GET}
      \item Endereço: \texttt{/api/alunos}
      \item Corpo: \textit{vazio} --- não precisa enviar dados
    \end{itemize}
  \end{block}

  \vspace{0.3em}
  \begin{block}{Resposta que o Spring devolve}
    \begin{itemize}
      \item Status: \texttt{200} --- deu certo
      \item Corpo: a lista em JSON
    \end{itemize}
  \end{block}
\end{frame}

\begin{frame}[fragile]{POST: envia dados no corpo}
  Exemplo: cadastrar a aluna Ana.

  \vspace{0.3em}
  \begin{block}{Pedido}
    \begin{itemize}
      \item Verbo: \texttt{POST}
      \item Endereço: \texttt{/api/alunos}
      \item Corpo:
    \end{itemize}
\begin{lstlisting}[basicstyle=\ttfamily\footnotesize]
{"nome": "Ana", "turma": "Entra21"}
\end{lstlisting}
  \end{block}

  \vspace{0.2em}
  \begin{block}{Resposta}
    Status \texttt{200} e o aluno já com \texttt{id}, em JSON.
  \end{block}
\end{frame}

\begin{frame}{O que é o corpo --- o body}
  \textbf{Corpo} = o recado que vai \textbf{dentro} do pedido,\\
  não na URL.

  \vspace{0.5em}
  \begin{itemize}
    \item Na Aula 02, o nome ia na URL:\\
          \texttt{/api/saudacao?nome=Ana}
    \item Hoje, no cadastro, nome e turma vão \textbf{dentro}\\
          do pedido, em JSON. Isso é o corpo.
  \end{itemize}

  \vspace{0.5em}
  A URL diz \textit{para quem} estamos falando
  \texttt{/api/alunos}.\\
  O corpo diz \textit{o que} queremos gravar.
\end{frame}

\begin{frame}{O Axios monta essa carta para a gente}
  Quando escrevemos no JavaScript:

  \vspace{0.4em}
  \begin{center}
    \texttt{axios.post("/api/alunos", \{ nome: "Ana", turma: "Entra21" \})}
  \end{center}

  \vspace{0.6em}
  O Axios faz, por baixo:
  \begin{enumerate}
    \item Verbo \texttt{POST}
    \item Endereço \texttt{/api/alunos}
    \item Corpo = aquele objeto, transformado em JSON
  \end{enumerate}

  \vspace{0.5em}
  No Java, \texttt{@RequestBody} significa:\\
  ``pega o JSON que veio no corpo e monta um objeto''.
\end{frame}

\begin{frame}{Como ver isso no navegador}
  Apertem \textbf{F12} $\rightarrow$ aba \textbf{Network} / Rede.

  \vspace{0.5em}
  Cliquem na linha do pedido \texttt{alunos}:
  \begin{itemize}
    \item \textbf{Headers} --- verbo e endereço
    \item \textbf{Payload} ou \textbf{Request} --- o corpo que
          \textbf{mandamos}
    \item \textbf{Response} --- o corpo que o Spring \textbf{devolveu}
  \end{itemize}

  \vspace{0.6em}
  Sempre que algo ``não gravar'', olhem essa aba\\
  antes de mudar o código de novo.
\end{frame}

% #############################################################################
\section{CRUD}

\begin{frame}{As quatro ações de um cadastro}
  \begin{center}
  \begin{tabular}{lll}
    \toprule
    Ação & Verbo HTTP & Exemplo \\
    \midrule
    Listar / ler & GET & Ver os alunos \\
    Cadastrar & POST & Incluir a Ana \\
    Editar & PUT & Trocar a turma da Ana \\
    Excluir & DELETE & Apagar a Ana \\
    \bottomrule
  \end{tabular}
  \end{center}

  \vspace{0.7em}
  CRUD é só o nome disso:\\
  Create, Read, Update, Delete.
\end{frame}

\begin{frame}{A lista precisa ser a mesma o tempo todo}
  Imaginem um caderno de chamada.

  \vspace{0.5em}
  \begin{itemize}
    \item GET = ler o caderno
    \item POST = escrever um nome novo \textbf{nesse} caderno
    \item PUT = corrigir um nome
    \item DELETE = apagar um nome
  \end{itemize}

  \vspace{0.6em}
  Se a cada leitura vocês pegam um caderno novo em branco,\\
  o nome que acabaram de escrever some.

  \vspace{0.4em}
  No código, o caderno é a \textbf{lista}.\\
  Ela precisa ficar \textbf{na classe do Controller},\\
  não dentro do método que lista.
\end{frame}

\begin{frame}[fragile]{Errado: lista dentro do método}
  Cada GET cria uma lista nova e joga fora no fim.

\begin{lstlisting}[language=Java,basicstyle=\ttfamily\footnotesize]
@GetMapping("/api/alunos")
public ArrayList<Aluno> listar() {
  ArrayList<Aluno> lista = new ArrayList<>();
  lista.add(new Aluno(1, "Ana", "Entra21"));
  return lista;
}
\end{lstlisting}

  \vspace{0.4em}
  O POST até pode adicionar um aluno\ldots\\
  mas no próximo GET essa lista já não existe.
\end{frame}

\begin{frame}[fragile]{Certo: lista na classe, métodos só usam}
  A lista continua \textbf{no Controller}.\\
  Só sai de \textbf{dentro da função} e vira campo da classe.

\begin{lstlisting}[language=Java,basicstyle=\ttfamily\footnotesize]
@RestController
public class AlunoController {

  private ArrayList<Aluno> alunos = new ArrayList<>();

  @GetMapping("/api/alunos")
  public ArrayList<Aluno> listar() {
    return alunos;
  }
}
\end{lstlisting}

  GET, POST, PUT e DELETE mexem nessa \textbf{mesma} lista.
\end{frame}

% #############################################################################
\section{Projeto cadastro-alunos}

\begin{frame}{Vamos criar um projeto novo}
  Primeira parte da aula: projeto \textbf{cadastro-alunos}.

  \vspace{0.5em}
  Serve para treinar CRUD com calma,\\
  e relembrar o \texttt{start.spring.io}.

  \vspace{0.6em}
  \begin{block}{O que esse projeto vai ter no final da Parte 1}
    Uma página com formulário e lista de alunos.\\
    Dá para cadastrar. Editar e excluir ficam nos desafios.
  \end{block}
\end{frame}

\begin{frame}{Passo 1 --- Gerar o projeto}
  Entrem em \texttt{https://start.spring.io}

  \vspace{0.4em}
  \begin{itemize}
    \item Project: Maven \quad Language: Java
    \item Group: \texttt{com.entra21}
    \item Artifact: \texttt{cadastro-alunos}
    \item Package name: \texttt{com.entra21.cadastro}
    \item Packaging: Jar \quad Java: 21
    \item Dependency: \textbf{Spring Web}
  \end{itemize}

  \vspace{0.4em}
  Generate $\rightarrow$ baixem $\rightarrow$ extraiam $\rightarrow$ abram a pasta
  no VS Code.
\end{frame}

\begin{frame}{Passo 2 --- Conferir se sobe}
  No terminal, dentro da pasta do projeto:

  \vspace{0.4em}
  \texttt{./mvnw spring-boot:run}\\
  {\small No Windows: \texttt{.\textbackslash mvnw.cmd spring-boot:run}}

  \vspace{0.5em}
  Se aparecer que o servidor está na porta 8080, está ok.\\
  Por enquanto não tem página nossa --- isso é o próximo passo.

  \vspace{0.5em}
  \begin{alertblock}{Dois Spring ao mesmo tempo}
    Se o StockSales ainda estiver rodando, a porta 8080
    já está ocupada. Parem o outro antes.
  \end{alertblock}
\end{frame}

\begin{frame}{Quais arquivos vamos criar}
  Projeto: \textbf{cadastro-alunos}

  \vspace{0.4em}
  \begin{center}
  \begin{tabular}{ll}
    \toprule
    Arquivo & Para quê \\
    \midrule
    \texttt{Aluno.java} & Os dados de um aluno \\
    \texttt{AlunoController.java} & A API: listar e cadastrar \\
    \texttt{static/index.html} & A página \\
    \texttt{static/js/app.js} & O JavaScript da página \\
    \bottomrule
  \end{tabular}
  \end{center}

  \vspace{0.5em}
  A classe \texttt{Application} o site já criou.\\
  Não precisamos mexer nela.
\end{frame}

\begin{frame}[fragile]{Onde cada arquivo fica}
\begin{lstlisting}[basicstyle=\ttfamily\footnotesize]
cadastro-alunos/
  src/main/java/com/entra21/cadastro/
    CadastroAlunosApplication.java   (ja existe)
    Aluno.java                       (vamos criar)
    AlunoController.java             (vamos criar)
  src/main/resources/static/
    index.html                       (vamos criar)
    js/app.js                        (vamos criar)
\end{lstlisting}

  Java fica em \texttt{java/}.\\
  Página e JS ficam em \texttt{resources/static/}.
\end{frame}

\begin{frame}[fragile]{Passo 3 --- Classe Aluno}
  \onde{Projeto \textbf{cadastro-alunos}\\
        Pasta \texttt{src/main/java/com/entra21/cadastro/}\\
        Arquivo novo: \texttt{Aluno.java}}

  Esta classe só guarda dados: id, nome e turma.\\
  Ainda não é a API.
\end{frame}

\begin{frame}[fragile]{Aluno.java --- atributos e construtores}
  \onde{\texttt{Aluno.java}}

\begin{lstlisting}[language=Java,basicstyle=\ttfamily\footnotesize]
package com.entra21.cadastro;

public class Aluno {
  private int id;
  private String nome;
  private String turma;

  public Aluno() {
  }

  public Aluno(int id, String nome, String turma) {
    this.id = id;
    this.nome = nome;
    this.turma = turma;
  }
}
\end{lstlisting}

  O construtor vazio é obrigatório para o POST funcionar.\\
  Completem os \texttt{get} e \texttt{set} dos três atributos.
\end{frame}

\begin{frame}{Por que dois construtores}
  \begin{itemize}
    \item \texttt{Aluno()} --- o Spring usa quando chega um JSON.
          Ele cria o objeto vazio e chama os \texttt{set}.
    \item \texttt{Aluno(id, nome, turma)} --- nós usamos para
          colocar Ana, Bruno e Carla na lista inicial.
  \end{itemize}

  \vspace{0.6em}
  Sem o construtor vazio, o cadastro costuma quebrar\\
  com erro no terminal.
\end{frame}

\begin{frame}[fragile]{Passo 4 --- Começar o Controller}
  \onde{Projeto \textbf{cadastro-alunos}\\
        Mesma pasta do \texttt{Aluno.java}\\
        Arquivo novo: \texttt{AlunoController.java}}

  Aqui fica a API. E aqui fica a lista ---\\
  \textbf{na classe}, não dentro do método \texttt{listar}.
\end{frame}

\begin{frame}[fragile]{AlunoController.java --- a lista}
  \onde{\texttt{AlunoController.java}}

\begin{lstlisting}[language=Java,basicstyle=\ttfamily\footnotesize]
package com.entra21.cadastro;

@RestController
public class AlunoController {

  private ArrayList<Aluno> alunos = new ArrayList<>();
  private int proximoId = 1;

  public AlunoController() {
    alunos.add(new Aluno(1, "Ana", "Entra21"));
    alunos.add(new Aluno(2, "Bruno", "Entra21"));
    alunos.add(new Aluno(3, "Carla", "Entra21"));
    proximoId = 4;
  }
}
\end{lstlisting}

  Não esqueçam os imports de \texttt{ArrayList}
  e \texttt{@RestController}.
\end{frame}

\begin{frame}[fragile]{GET --- só devolve a lista da classe}
  \onde{Ainda em \texttt{AlunoController.java},\\
        um método novo dentro da classe}

\begin{lstlisting}[language=Java,basicstyle=\ttfamily\footnotesize]
@GetMapping("/api/alunos")
public ArrayList<Aluno> listar() {
  return alunos;
}
\end{lstlisting}

  \vspace{0.4em}
  Salvem, reiniciem o Spring, abram no navegador:

  \vspace{0.3em}
  \texttt{http://localhost:8080/api/alunos}

  \vspace{0.4em}
  Tem que aparecer o JSON de Ana, Bruno e Carla.\\
  Ainda sem página --- só testando a API.
\end{frame}

\begin{frame}[fragile]{POST --- lê o corpo e adiciona na lista}
  \onde{Ainda em \texttt{AlunoController.java}}

\begin{lstlisting}[language=Java,basicstyle=\ttfamily\footnotesize]
@PostMapping("/api/alunos")
public Aluno cadastrar(@RequestBody Aluno aluno) {
  aluno.setId(proximoId);
  proximoId = proximoId + 1;
  alunos.add(aluno);
  return aluno;
}
\end{lstlisting}

  \vspace{0.3em}
  \texttt{@RequestBody} = pega o JSON do \textbf{corpo} do pedido.

  \vspace{0.3em}
  Imports novos: \texttt{PostMapping} e \texttt{RequestBody}.
\end{frame}

\begin{frame}{O que o POST faz, em ordem}
  \begin{enumerate}
    \item O JavaScript manda JSON no corpo:\\
          nome e turma
    \item O Spring monta um \texttt{Aluno} com os \texttt{set}
    \item Nós colocamos o próximo \texttt{id}
    \item Adicionamos na lista \textbf{da classe}
    \item Devolvemos o aluno já com id
  \end{enumerate}

  \vspace{0.5em}
  O GET seguinte vê o aluno novo,\\
  porque a lista não foi recriada.
\end{frame}

\begin{frame}[fragile]{Passo 5 --- A página}
  \onde{Projeto \textbf{cadastro-alunos}\\
        Pasta \texttt{src/main/resources/static/}\\
        Arquivo novo: \texttt{index.html}}

  Esta é a tela. Ainda não tem lógica --- só o desenho:\\
  um formulário e uma lista vazia.
\end{frame}

\begin{frame}[fragile]{index.html}
  \onde{\texttt{src/main/resources/static/index.html}}

\begin{lstlisting}[basicstyle=\ttfamily\footnotesize]
<h1>Cadastro de alunos</h1>

<form id="form-aluno">
  <input id="nome" type="text" placeholder="Nome">
  <input id="turma" type="text" placeholder="Turma">
  <button type="submit">Cadastrar</button>
</form>

<ul id="lista-alunos"></ul>

<script src="https://cdn.jsdelivr.net/npm/axios/dist/axios.min.js"></script>
<script src="/js/app.js"></script>
\end{lstlisting}

  Axios \textbf{antes} do \texttt{app.js}.
\end{frame}

\begin{frame}{Passo 6 --- O JavaScript}
  \onde{Projeto \textbf{cadastro-alunos}\\
        Pasta \texttt{src/main/resources/static/js/}\\
        Arquivo novo: \texttt{app.js}}

  Duas tarefas neste arquivo:
  \begin{enumerate}
    \item Ao abrir a página, buscar a lista na API
    \item Ao enviar o formulário, cadastrar na API\\
          e atualizar a lista
  \end{enumerate}
\end{frame}

\begin{frame}[fragile]{app.js --- listar}
  \onde{\texttt{static/js/app.js}}

\begin{lstlisting}[basicstyle=\ttfamily\footnotesize]
var lista = document.getElementById("lista-alunos");

function carregarAlunos() {
  axios.get("/api/alunos")
    .then(function (resposta) {
      var alunos = resposta.data;
      lista.innerHTML = "";
      for (var i = 0; i < alunos.length; i++) {
        var li = document.createElement("li");
        li.textContent =
          alunos[i].nome + " - " + alunos[i].turma;
        lista.appendChild(li);
      }
    });
}

carregarAlunos();
\end{lstlisting}
\end{frame}

\begin{frame}{Por que uma função carregarAlunos}
  Vamos chamar essa função em dois momentos:

  \vspace{0.5em}
  \begin{itemize}
    \item Quando a página abre
    \item Depois de cadastrar --- para a lista mostrar o novo
  \end{itemize}

  \vspace{0.6em}
  Se cadastrar e não chamar de novo, o backend gravou,\\
  mas a tela fica antiga.
\end{frame}

\begin{frame}[fragile]{app.js --- cadastrar}
  \onde{Ainda em \texttt{static/js/app.js}, abaixo do listar}

  O formulário, sozinho, recarrega a página.\\
  \texttt{preventDefault} impede isso.

\begin{lstlisting}[basicstyle=\ttfamily\footnotesize]
var form = document.getElementById("form-aluno");

form.addEventListener("submit", function (evento) {
  evento.preventDefault();
  var nome = document.getElementById("nome").value;
  var turma = document.getElementById("turma").value;

  axios.post("/api/alunos", { nome: nome, turma: turma })
    .then(function () {
      carregarAlunos();
    });
});
\end{lstlisting}
\end{frame}

\begin{frame}{O caminho do cadastro, do clique até a tela}
  \begin{enumerate}
    \item A pessoa preenche nome e turma e clica em Cadastrar
    \item O JS manda \texttt{POST /api/alunos} com JSON no corpo
    \item O Controller lê o corpo, gera o id, coloca na lista
    \item O JS recebe a resposta e chama \texttt{carregarAlunos}
    \item O GET busca a lista de novo e redesenha a \texttt{<ul>}
  \end{enumerate}

  \vspace{0.5em}
  No F12, aba Network, vocês devem ver\\
  um POST e, em seguida, um GET.
\end{frame}

\begin{frame}{Conferindo o exemplo juntos}
  \begin{itemize}
    \item [ ] \texttt{http://localhost:8080/} abre o formulário
    \item [ ] Ana, Bruno e Carla aparecem
    \item [ ] Cadastrar um nome novo atualiza a lista
    \item [ ] Recarregar a página \textbf{mantém} o nome novo
    \item [ ] Parar e subir o Spring de novo: voltam só os três
  \end{itemize}

  \vspace{0.5em}
  Se o nome some ao recarregar a página,\\
  a lista ainda está sendo criada dentro do \texttt{listar}.
\end{frame}

\begin{frame}{Erros que mais aparecem}
  \begin{itemize}
    \item Porta 8080 ocupada --- outro Spring ainda rodando
    \item \texttt{axios is not defined} --- script do Axios
          depois do \texttt{app.js}
    \item POST quebra --- faltou construtor vazio no
          \texttt{Aluno}
    \item Página recarrega e o form esvazia --- faltou
          \texttt{preventDefault}
    \item Cadastrou e a lista não muda --- não chamou
          \texttt{carregarAlunos}
  \end{itemize}
\end{frame}

% #############################################################################
\section{Desafios}

\begin{frame}{Agora é com vocês}
  Continuem no projeto \textbf{cadastro-alunos}.\\
  Mesmos arquivos. Não criem outro projeto.

  \vspace{0.5em}
  \begin{enumerate}
    \item Excluir aluno
    \item Editar no backend
    \item Editar pela tela
  \end{enumerate}

  \vspace{0.5em}
  Salve $\rightarrow$ reinicie o Spring $\rightarrow$ teste no navegador.\\
  F12, aba Network.
\end{frame}

\begin{frame}{PUT e DELETE usam o id no endereço}
  Para mudar \textbf{um} aluno, o servidor precisa saber
  \textbf{qual}.

  \vspace{0.5em}
  \begin{center}
  \begin{tabular}{ll}
    \toprule
    Ação & Pedido \\
    \midrule
    Editar o aluno 2 & \texttt{PUT /api/alunos/2} \\
      & corpo com nome e turma novos \\
    Excluir o aluno 2 & \texttt{DELETE /api/alunos/2} \\
      & sem corpo \\
    \bottomrule
  \end{tabular}
  \end{center}

  \vspace{0.5em}
  O \texttt{2} é o \texttt{@PathVariable}.
\end{frame}

\begin{frame}{Desafio 1 --- Excluir}
  \onde{1) \texttt{AlunoController.java}\\
        2) \texttt{static/js/app.js}}

  No Controller, método novo:

  \vspace{0.3em}
  \texttt{DELETE /api/alunos/\{id\}}

  \vspace{0.4em}
  Percorram a lista. Se o id bater:
  \texttt{alunos.remove(i)} e parem o \texttt{for}.

  \vspace{0.4em}
  No JS, cada item da lista ganha um botão Excluir.\\
  No clique: \texttt{axios.delete("/api/alunos/" + id)}\\
  e depois \texttt{carregarAlunos()}.
\end{frame}

\begin{frame}[fragile]{Desafio 1 --- o botão no JS}
  \onde{\texttt{static/js/app.js}, dentro do \texttt{for}\\
        que monta a lista}

\begin{lstlisting}[basicstyle=\ttfamily\footnotesize]
var botao = document.createElement("button");
botao.textContent = "Excluir";
botao.setAttribute("data-id", alunos[i].id);

botao.addEventListener("click", function (evento) {
  var id = evento.target.getAttribute("data-id");
  axios.delete("/api/alunos/" + id)
    .then(function () {
      carregarAlunos();
    });
});

li.appendChild(botao);
\end{lstlisting}

  Usem \texttt{evento.target}, não a variável \texttt{botao}.
\end{frame}

\begin{frame}[fragile]{Desafio 2 --- Editar no Controller}
  \onde{\texttt{AlunoController.java}}

  Método novo: \texttt{PUT /api/alunos/\{id\}}

  \vspace{0.3em}
  Recebam o id da URL e o aluno do corpo.\\
  Achem na lista, atualizem nome e turma, devolvam o aluno.

\begin{lstlisting}[language=Java,basicstyle=\ttfamily\footnotesize]
@PutMapping("/api/alunos/{id}")
public Aluno atualizar(@PathVariable int id,
                       @RequestBody Aluno dados) {
  // for na lista, se o id bater:
  // setNome e setTurma, depois return
  return null;
}
\end{lstlisting}
\end{frame}

\begin{frame}{Desafio 3 --- Editar pela tela}
  \onde{\texttt{static/js/app.js}}

  Objetivo: clicar em Editar, o formulário se preenche,\\
  e o Cadastrar passa a salvar a alteração.

  \vspace{0.4em}
  \begin{enumerate}
    \item Variável \texttt{idEditando}, começando vazia
    \item Botão Editar coloca o id, o nome e a turma no form
    \item Se \texttt{idEditando} estiver vazio $\rightarrow$ POST
    \item Se tiver id $\rightarrow$ PUT para aquele id
    \item No fim, limpem o form e o \texttt{idEditando}
  \end{enumerate}
\end{frame}

\begin{frame}{Extra}
  Se terminarem os três, ainda no \textbf{cadastro-alunos}:

  \vspace{0.4em}
  \begin{itemize}
    \item Não cadastrar se o nome estiver vazio
    \item Mostrar o id na lista
    \item Texto na página: ``Aluno cadastrado''
  \end{itemize}
\end{frame}

\begin{frame}{Checklist}
  \begin{itemize}
    \item [ ] Estou no projeto \texttt{cadastro-alunos}?
    \item [ ] A lista é campo da classe, não do método
              \texttt{listar}?
    \item [ ] \texttt{Aluno} tem construtor vazio?
    \item [ ] POST usa \texttt{@RequestBody}?
    \item [ ] O form tem \texttt{preventDefault}?
    \item [ ] Depois de gravar, chama \texttt{carregarAlunos}?
  \end{itemize}
\end{frame}

\begin{frame}{Intervalo}
  \begin{center}
    \Large Retomamos às \textbf{20h20}\\[0.8em]
    \normalsize
    Fechem o \texttt{cadastro-alunos}.\\
    Na segunda parte abrimos o \textbf{StockSales}.
  \end{center}
\end{frame}

% #############################################################################
\section{Parte 2 --- StockSales}

\begin{frame}{Mudamos de projeto}
  \onde{Agora é o projeto \textbf{StockSales},\\
        não o cadastro-alunos.}

  Parem o Spring do cadastro-alunos.\\
  Abram a pasta do StockSales e subam de novo.
\end{frame}

\begin{frame}{Onde o StockSales parou}
  Arquivos que vocês já têm:

  \vspace{0.4em}
  \begin{itemize}
    \item \texttt{controller/ProdutoController.java}
    \item \texttt{model/Produto.java}
    \item \texttt{static/produtos.html}
    \item \texttt{static/js/produtos.js}
  \end{itemize}

  \vspace{0.5em}
  Hoje a lista é criada \textbf{dentro} de
  \texttt{listarProdutos}.\\
  Por isso todo GET volta sempre os quatro produtos fixos.

  \vspace{0.4em}
  A lista deve \textbf{continuar no Controller}.\\
  Só precisa sair de dentro dessa função.
\end{frame}

\begin{frame}{O que vamos fazer no StockSales}
  \begin{enumerate}
    \item Subir a lista de produtos para a classe
          \texttt{ProdutoController}
    \item Cadastrar, editar e excluir produto
    \item Formulário na tela de produtos
    \item Criar cliente: classe, Controller e tela
  \end{enumerate}

  \vspace{0.5em}
  É o mesmo raciocínio do cadastro-alunos.\\
  Só mudam os campos e os nomes dos arquivos.
\end{frame}

\begin{frame}[fragile]{Produto.java --- construtor vazio}
  \onde{Projeto \textbf{StockSales}\\
        Arquivo \texttt{src/main/java/com/entra21/stocksales/model/Produto.java}}

  Acrescentem este construtor. Mantenham o outro.

\begin{lstlisting}[language=Java,basicstyle=\ttfamily\footnotesize]
public Produto() {
}
\end{lstlisting}

  Sem ele, o POST de produto quebra.
\end{frame}

\begin{frame}[fragile]{A lista sai do método, fica na classe}
  \onde{\texttt{controller/ProdutoController.java}}

  \textbf{Não} tirem a lista do Controller.\\
  Tirem só de dentro de \texttt{listarProdutos}.

\begin{lstlisting}[language=Java,basicstyle=\ttfamily\footnotesize]
private List<Produto> produtos = new ArrayList<>();
private int proximoId = 1;

public ProdutoController() {
  produtos.add(new Produto(1, "Teclado Mecanico", 250, 12));
  produtos.add(new Produto(2, "Mouse Gamer", 120, 30));
  produtos.add(new Produto(3, "Monitor 24", 900, 8));
  produtos.add(new Produto(4, "Headset", 200, 15));
  proximoId = 5;
}
\end{lstlisting}
\end{frame}

\begin{frame}[fragile]{listarProdutos fica só isto}
  \onde{Ainda em \texttt{ProdutoController.java}}

  Antes o método criava a lista. Agora só devolve.

\begin{lstlisting}[language=Java,basicstyle=\ttfamily\footnotesize]
@GetMapping("/api/produtos")
public List<Produto> listarProdutos() {
  return produtos;
}
\end{lstlisting}

  Testem \texttt{/api/produtos} no navegador:\\
  tem que continuar mostrando os quatro.
\end{frame}

\begin{frame}[fragile]{POST de produto}
  \onde{Ainda em \texttt{ProdutoController.java}}

\begin{lstlisting}[language=Java,basicstyle=\ttfamily\footnotesize]
@PostMapping("/api/produtos")
public Produto cadastrar(@RequestBody Produto produto) {
  produto.setId(proximoId);
  proximoId = proximoId + 1;
  produtos.add(produto);
  return produto;
}
\end{lstlisting}

  PUT e DELETE: mesmo \texttt{for} do desafio dos alunos,\\
  com \texttt{/api/produtos/\{id\}}.
\end{frame}

\begin{frame}{Tela de produtos}
  \onde{\texttt{static/produtos.html} e \texttt{static/js/produtos.js}}

  Na HTML, acima da lista: formulário com nome, preço e estoque.

  \vspace{0.4em}
  No JS:
  \begin{itemize}
    \item Função \texttt{carregarProdutos} --- vocês já listam;
          só extraíam para uma função
    \item \texttt{preventDefault} no form
    \item \texttt{axios.post("/api/produtos", \{ ... \})}
    \item Botão Excluir em cada item
  \end{itemize}
\end{frame}

\begin{frame}[fragile]{Preço e estoque são número}
  \onde{\texttt{static/js/produtos.js}}

  O input sempre devolve texto. O Java espera número.

\begin{lstlisting}[basicstyle=\ttfamily\footnotesize]
axios.post("/api/produtos", {
  nome: document.getElementById("nome").value,
  preco: Number(document.getElementById("preco").value),
  estoque: Number(document.getElementById("estoque").value)
});
\end{lstlisting}
\end{frame}

\begin{frame}[fragile]{Classe Cliente}
  \onde{Projeto \textbf{StockSales}\\
        Arquivo novo:\\
        \texttt{src/main/java/com/entra21/stocksales/model/Cliente.java}}

\begin{lstlisting}[language=Java,basicstyle=\ttfamily\footnotesize]
public class Cliente {
  private int id;
  private String nome;
  private String email;
  private String telefone;

  public Cliente() {}
}
\end{lstlisting}

  Construtor com os quatro valores + gets e sets.
\end{frame}

\begin{frame}{Controller e tela de clientes}
  \onde{Arquivos novos:\\
        \texttt{controller/ClienteController.java}\\
        \texttt{static/clientes.html}\\
        \texttt{static/js/clientes.js}}

  Copiem o esquema do produto:
  \begin{itemize}
    \item Lista na \textbf{classe} do Controller
    \item GET / POST / PUT / DELETE em \texttt{/api/clientes}
    \item Formulário: nome, e-mail, telefone
  \end{itemize}

  \vspace{0.4em}
  Na home, um link para \texttt{/clientes.html}.
\end{frame}

\begin{frame}{O que não entra hoje}
  \begin{itemize}
    \item Vendas e baixa de estoque --- próxima aula
    \item Banco de dados
    \item Login
  \end{itemize}

  \vspace{0.5em}
  \begin{block}{Meta no StockSales}
    Cadastrar, listar, editar e excluir produtos.\\
    Cadastrar, listar, editar e excluir clientes.\\
    Lista na classe do Controller, ainda sem banco.
  \end{block}
\end{frame}

\begin{frame}{Critérios de aceite}
  \begin{itemize}
    \item [ ] Produto novo aparece na lista
    \item [ ] Recarregar a página mantém, até reiniciar o Spring
    \item [ ] Excluir some da lista
    \item [ ] Editar muda nome, preço ou estoque
    \item [ ] Clientes no mesmo esquema
    \item [ ] Home leva a produtos e a clientes
  \end{itemize}
\end{frame}

% #############################################################################
\section{Recapitulando}

\begin{frame}{O que vimos hoje}
  \begin{enumerate}
    \item Requisição: verbo + endereço + corpo
    \item Resposta: status + JSON
    \item POST/PUT mandam dados no corpo
    \item A lista fica no Controller, fora do método de listar
    \item CRUD no cadastro-alunos e no StockSales
  \end{enumerate}
\end{frame}

\begin{frame}{Na próxima aula}
  \begin{block}{Aula 04}
    \begin{itemize}
      \item Vendas no StockSales
      \item Escolher cliente e produtos
      \item Baixa de estoque
    \end{itemize}
  \end{block}

  \vspace{0.6em}
  Tragam o StockSales com produtos e clientes cadastráveis.
\end{frame}

\begin{frame}
  \begin{center}
    \Huge Obrigado!\\[1.2em]
    \normalsize Mauricio Duailibi Neto\\
    Turma Java Entra21
  \end{center}
\end{frame}

\end{document}
